\chapter{The ORTi3D Graphical User Interface}
\def\shorttitle{ORTi3D} %

\section{Before you start modelling}

ORTi is an open source graphical user interface that provides a modelling environment for multiple reactive transport models, including the MODFLOW/MT3DMS-based model PHT3D.  
ORTi can be used in two different interface modes. The first mode is a windows/python based mode, termed �classical� mode, which has been developed and used over the last few years. The second and just recently introduced mode is the QGIS mode. 
In order to allow constructed models to remain interchangeable between both modes the libraries that are required 
by the two modes need to be stored in a single QGIS plugin folder. This will allow any ORTi model to be saved for subsequent processing in either QQIS or classical interface mode.

To start working in the classical interface mode (after installing the QGIS version), just double click on:
C:\Users\�yourusername�\AppData\Roaming\QGIS\QGIS3\profiles\default\python\plugins\qORTi3d\startOrti3d.bat

\section{Installing the plugin}
Go to \\
https://www.orti3d.org/documents/ORTi3D_qgis_plugin_version.zip \\

to download the QGIS plugin. You have to unzip it at a specific location: 

C:\Users\�your username�\AppData\Roaming\QGIS\QGIS3\profiles\default\python\plugins\qORTi3d \\

You need to create the folder qORTi3d in the plugin folder. 
Once the file is unzipped verify that the ilibq folder is just below the qOrti3d folder.
Once QGIS is opened, go to the Plugins menu and select manage plugins, 
find the qORTi3d plugin and select it. You should now see it at the bottom of the Plugins menu.
To install the lastest version, open the plugin, go in the file > Help and click on download last. 
You need to restart QGIS (later you can install a plugin that is called plugin reloader that allows you to reload a plugin that has been modified without restarting QGIS, this is very useful for developers)
